\section{RQ7: What Are the Principal Barriers to Syntactic Correctness in Rocq?}
\label{appendix:rq7}

While state-of-the-art models achieve relatively high syntax compliance, a non-negligible fraction of generated specifications still fails to compile (as reflected in the \textsc{Syntax} metric in Table~\ref{tab:generation}). To understand the root causes, we conducted a qualitative analysis of all 156 specifications that failed the \textsc{Syntax} stage across the top three performing models: Claude 4.5 Opus, GPT-5, and Gemini 3 Pro Preview. We classified the errors into three primary categories:

\begin{itemize}
    \item \textbf{Type Mismatches (45\%):} Logical errors where the model conflates distinct Rocq types, such as Peano natural numbers (\texttt{nat}), binary integers (\texttt{Z}), and propositions (\texttt{Prop}).
    \item \textbf{Missing Dependencies (22\%):} Cases where the model invokes semantically appropriate functions (e.g., from \texttt{Ascii} or \texttt{ZArith}) but fails to include necessary \texttt{Require Import} statements.
    \item \textbf{Ill-formed Recursion (17\%):} \texttt{Fixpoint} definitions that violate the \textit{Calculus of Inductive Constructions}' strict structural termination requirements, often by recursing on non-syntactic sub-terms.
\end{itemize}

\paragraph{Error Distribution and Analysis.}
Table~\ref{tab:error-analysis} summarizes the error distribution. These failures reflect a fundamental tension between the probabilistic nature of LLMs and the rigid logical kernels of proof assistants.

\begin{table}[h]
\centering
\caption{Breakdown of compilation errors (N=156) by category.}
\label{tab:error-analysis}
\vskip 0.1in
\begin{tabular}{lccccc}
\toprule
\textbf{Model} & \textbf{Type Mismatch} & \textbf{Missing Dep.} & \textbf{Bad Recursion} & \textbf{Others} & \textbf{Total} \\
\midrule
Claude 4.5 Opus & 24 & 19 & 12 & 11 & 66 \\
GPT-5           & 29 & 7  & 9  & 9  & 54 \\
Gemini 3 Pro Preview   & 17 & 8  & 5  & 6  & \textbf{36} \\
\midrule
\textbf{Total}  & \textbf{70} (45\%) & \textbf{34} (22\%) & \textbf{26} (17\%) & \textbf{26} (17\%) & 156 \\
\bottomrule
\end{tabular}
\end{table}

\textbf{Type System Rigidity vs. Language Bias.}
Type Mismatches remain the dominant failure. Models frequently treat \texttt{nat} and \texttt{Z} interchangeably, likely an artifact of pre-training on dynamically typed languages like Python. Unlike those languages, Rocq requires explicit casting (e.g., \texttt{Z.of\_nat}) due to the fundamentally different inductive structures of these types.

\textbf{Context Management and Imports.}
Claude 4.5 Opus exhibited a notably higher rate of Missing Dependencies (19 cases). This suggests a "knowledge-action gap": the model understands the semantic need for a library function but neglects the syntactic boilerplate required to make it available in the current environment. Gemini 3 Pro Preview demonstrated more robust context awareness in this regard.

\textbf{The Structural Termination Gap.}
All models struggle with Ill-formed Recursive Definitions (17\% of errors). Rocq's guard checker requires recursive functions to be structurally terminating, meaning recursive calls must be performed on direct syntactic sub-terms of the inductive input. We observed that models often generate ``Pythonic'' recursion patterns that are mathematically sound but formally rejected. 

In this section, we provide a concrete example of the \textbf{Ill-formed Recursion} failure mode, which accounts for approximately 17\% of compilation errors. Figure~\ref{fig:bad_recursion_case} illustrates a case where the model attempts to generate a specification for a sequence similar to Fibonacci (specifically, the \texttt{fib4} sequence).

\begin{figure}[h]
    \centering
    % colback=blue!3!white : 极淡的蓝色背景
    % colframe=blue!60!black : 深蓝色边框
    \begin{tcolorbox}[
        enhanced,
        title=\textbf{Case Study: Ill-formed Recursion in \texttt{fib4\_seq}},
        fonttitle=\bfseries\small,
        colback=blue!3!white,   
        colframe=blue!50!black,    
        boxrule=0.8pt,
        arc=2pt,
        breakable
    ]
    
    % --- 1. The Generated Code ---
    \textbf{Generated Specification:}
    \begin{lstlisting}[
        language=Coq, 
        frame=single, 
        framerule=0pt,             
        backgroundcolor=\color{white}, 
        basicstyle=\ttfamily\small
    ]
Require Import Coq.ZArith.ZArith.
Require Import Coq.Arith.PeanoNat.

Open Scope Z_scope.

Fixpoint fib4_seq (n : nat) : Z :=
  match n with
  | 0%nat => 0
  | 1%nat => 0
  | 2%nat => 2
  | 3%nat => 0
  | S (S (S (S m))) => fib4_seq (S (S (S m))) + fib4_seq (S (S m)) + 
                       fib4_seq (S m) + fib4_seq m
  end.

Definition fib4_spec (n : Z) (res : Z) : Prop :=
  n >= 0 /\ res = fib4_seq (Z.to_nat n).
    \end{lstlisting}

    \vspace{0.2cm}
    
    % --- 2. The Compiler Error ---
    \textbf{Rocq Compiler Error:}
    \begin{tcolorbox}[
        colback=red!5!white,
        colframe=red!75!black,
        left=2pt, right=2pt, top=2pt, bottom=2pt,
        arc=0mm, boxrule=0.5pt,
        title=Compiler Output,
        fonttitle=\bfseries\scriptsize,
        coltitle=white
    ]
    \footnotesize\ttfamily\color{red!40!black}
    Recursive definition of fib4\_seq is ill-formed.\\
    In environment\\
    fib4\_seq : nat -> Z\\
    ...\\
    Recursive call to fib4\_seq has principal argument equal to "S m" instead of one of the following variables: "n0" "n1" "n2" "m".
    \end{tcolorbox}
    
    \end{tcolorbox}
    
    \caption{\textbf{Example of Ill-formed Recursion.} The model attempts to implement a recurrence relation $f(n) = \sum_{i=1}^{4} f(n-i)$ but fails Rocq's structural termination check. Rocq requires recursive arguments to be strict subterms of the matched pattern (e.g., \texttt{m}), but the model passes \texttt{S m} (conceptually $n-3$), which Rocq cannot automatically verify as terminating without additional measures.}
    \label{fig:bad_recursion_case}
\end{figure}